\documentclass[10pt,aspectratio=169]{beamer}

\usepackage[utf8]{inputenc}
\usepackage[T1]{fontenc}
\usepackage[french]{babel}
\usepackage{amsmath,amssymb,amsfonts}
\usepackage{booktabs}
\usepackage{tikz}
\usepackage{graphicx}
\usepackage{array}
\usepackage{colortbl}
\usepackage{lmodern}
\usepackage{tcolorbox}
\usepackage{pifont}
\usepackage{csquotes}
\usetikzlibrary{positioning,calc,arrows.meta,shapes.geometric,fit}
\tcbuselibrary{skins}

\definecolor{primary}{RGB}{12,54,106}
\definecolor{secondary}{RGB}{234,241,255}
\definecolor{accent}{RGB}{255,111,0}
\definecolor{darktext}{RGB}{20,33,61}
\definecolor{lighttext}{RGB}{108,117,125}
\definecolor{success}{RGB}{42,157,143}
\definecolor{danger}{RGB}{230,57,70}
\definecolor{warning}{RGB}{244,162,97}

\usefonttheme{professionalfonts}
\setbeamerfont{title}{size=\huge,series=\bfseries}
\setbeamerfont{subtitle}{size=\large,series=\normalfont}
\setbeamerfont{frametitle}{size=\Large,series=\bfseries}
\setbeamercolor{background canvas}{bg=white}
\setbeamercolor{normal text}{fg=darktext}
\setbeamercolor{frametitle}{fg=primary}
\setbeamercolor{title}{fg=primary}
\setbeamercolor{subtitle}{fg=lighttext}
\setbeamercolor{structure}{fg=primary}
\setbeamercolor{item}{fg=accent}

\newtcolorbox{keybox}[1][] {
    colframe=primary!25,
    colback=secondary!35,
    boxrule=0.8pt,
    arc=4pt,
    left=6pt,
    right=6pt,
    top=5pt,
    bottom=5pt,
    #1
}

\newcommand{\cmark}{\textcolor{success}{\ding{51}}}
\newcommand{\xmark}{\textcolor{danger}{\ding{55}}}

\setbeamertemplate{navigation symbols}{}

\begin{document}

\setbeamertemplate{footline}{}
\begin{frame}[plain]
\begin{tikzpicture}[remember picture,overlay]
    \fill[secondary!30] (current page.south west) rectangle (current page.north east);
    \fill[primary!5, opacity=0.85] (current page.north west) -- (current page.south east) -- (current page.south west) -- cycle;
    \fill[accent, opacity=0.08] ($(current page.center)+(3cm,1.2cm)$) circle (4cm);
    \node[anchor=center, align=center] at (current page.center) {
        {\Huge\bfseries\textcolor{primary}{SpeechCoach}}\\[0.3cm]
        {\huge\bfseries\textcolor{darktext}{Réunion d'avancement n\textdegree 4}}\\[0.55cm]
        \begin{tikzpicture}
            \fill[accent] (0,0) rectangle (3.2cm,0.12cm);
        \end{tikzpicture}\\[0.55cm]
        {\Large\itshape\textcolor{lighttext}{Optimisation ASR multilingue, analyse visuelle adaptative}}\\[0.2cm]
{\Large\itshape\textcolor{lighttext}{et coaching LLM hybride}}\\[1.2cm]
        {\normalsize Hicham Moussaid}\\[0.15cm]
        {\small Master SIE -- ENS Meknès}\\[0.15cm]
        {\small Encadrant : Pr. Abdelaaziz Hessane}
    };
\end{tikzpicture}
\end{frame}

\setbeamertemplate{footline}{
    \begin{tikzpicture}[remember picture,overlay]
        \fill[primary!5] (current page.south west) rectangle ($(current page.south east)+(0,15pt)$);
        \node[anchor=south east, xshift=-14pt, yshift=4pt, text=primary!60, font=\scriptsize\bfseries] at (current page.south east) {
            \insertframenumber~|~\inserttotalframenumber
        };
        \fill[accent] (current page.south west) rectangle ($(current page.south west)+(3pt,15pt)$);
        \node[anchor=south west, xshift=10pt, yshift=4pt, text=primary!60, font=\tiny\bfseries] at (current page.south west) {
            SPEECHCOACH | AVANCEMENT 2026
        };
    \end{tikzpicture}
}

\begin{frame}{Plan de la présentation}
\begin{columns}[T]
\begin{column}{0.48\textwidth}
\textbf{1. Objectifs et contexte}
\begin{itemize}
    \item trois axes de travail
    \item évolution du coaching
\end{itemize}
\vspace{0.3cm}
\textbf{2. Travail ASR}
\begin{itemize}
    \item problème et benchmark
    \item correctifs et résultats
\end{itemize}
\end{column}
\begin{column}{0.48\textwidth}
\textbf{3. Analyse visuelle adaptative}
\begin{itemize}
    \item problème smartphone
    \item profils et auto-détection
\end{itemize}
\vspace{0.3cm}
\textbf{4. Stratégie LLM et architecture}
\begin{itemize}
    \item fine-tuning et transition Groq
    \item performance et prochaines étapes
\end{itemize}
\end{column}
\end{columns}
\end{frame}

\begin{frame}{Objectifs et évolution du coaching}
\begin{columns}[T]
\begin{column}{0.49\textwidth}
\begin{tcolorbox}[title=Objectifs de la phase, colframe=primary, colback=primary!3]
\small
\begin{enumerate}
    \item Améliorer le bloc d'exercices
    \item Optimiser ASR multilingue (fillers, répétitions)
    \item Renforcer cohérence moteur--LLM
\end{enumerate}
\end{tcolorbox}
\end{column}
\begin{column}{0.49\textwidth}
\begin{tcolorbox}[title=Évolution du coaching, colframe=accent, colback=accent!3]
\small
\textbf{Phase 1 :} Programme fixe 7 jours \xmark\\
\textit{Problème :} utilisateur pas prêt à s'engager sur une semaine\\
\textbf{Phase 2 :} Approche par modes \xmark\\
\textit{Problème :} templates limitent créativité IA\\
\textbf{Phase 3 :} Master Coach Agent \cmark\\
\textit{Solution :} génération dynamique sans templates
\end{tcolorbox}
\end{column}
\end{columns}
\vspace{0.3cm}
\begin{keybox}
\textbf{Master Coach Agent :} IA reçoit métriques actuelles + RAG + \textbf{historique des 3 dernières sessions complétées} (scores par axe : voix, présence, gestes, cadrage ; tendances récentes ; chronologie détaillée), conçoit mission unique en 3 étapes adaptée au parcours utilisateur, post-processeur assure focus sur un seul axe
\end{keybox}
\end{frame}

\begin{frame}{Master Coach Agent : Détails de l'implémentation}
\begin{columns}[T]
\begin{column}{0.49\textwidth}
\begin{tcolorbox}[title=Flux de génération, colframe=primary, colback=primary!3]
\small
\textbf{Phase immédiate -- scoring :}
\begin{itemize}
    \item Moteur déterministe calcule métriques brutes
    \item Tableau de bord s'actualise sans attendre LLM
\end{itemize}
\textbf{Phase différée -- enrichissement :}
\begin{itemize}
    \item Worker Celery agrège historique + RAG
    \item Appel API Groq pour coaching
\end{itemize}
\end{tcolorbox}
\end{column}
\begin{column}{0.49\textwidth}
\begin{tcolorbox}[title=Intelligence pédagogique, colframe=success, colback=success!3]
\small
\textbf{RAG :}
\begin{itemize}
    \item Base documentaire via embeddings
    \item Recherche vectorielle FAISS
    \item Sans recopie littérale des fiches
\end{itemize}
\textbf{Sortie JSON structurée :}
\begin{itemize}
    \item Bilan global, point prioritaire
    \item Encouragement
    \item Exercice (titre + 3 étapes)
    \item Axe de focus (voix, présence, gestuelle, scène)
\end{itemize}
\end{tcolorbox}
\end{column}
\end{columns}
\end{frame}

\begin{frame}{Travail ASR : Problème et benchmark}
\begin{columns}[T]
\begin{column}{0.49\textwidth}
\begin{tcolorbox}[title=Problème identifié, colframe=danger, colback=danger!3]
\small
\begin{itemize}
    \item Fillers (\textit{um}, \textit{uh}, \textit{euh}, \textit{mmm}, arabes) entendus mais non conservés
    \item Affecte directement le score voix/débit
    \item Incohérence entre analyse et transcription
\end{itemize}
\end{tcolorbox}
\end{column}
\begin{column}{0.49\textwidth}
\begin{tcolorbox}[title=Benchmark CrisperWhisper, colframe=warning!60!black, colback=warning!8]
\small
\textbf{Résultats :}
\begin{itemize}
    \item \xmark~Temps d'exécution trop élevé
    \item \xmark~Amélioration non convaincante
    \item \xmark~Qualité insuffisante
\end{itemize}
\textbf{Conclusion :} rejeté, conservé Faster-Whisper
\end{tcolorbox}
\end{column}
\end{columns}
\vspace{0.3cm}
\begin{keybox}
\textbf{5 sources d'erreur identifiées :} chaînes arabes corrompues, prompt multilingue, listes incomplètes, lissage sons nasaux, modélisation répétitions insuffisante
\end{keybox}
\end{frame}

\begin{frame}{Travail ASR : Correctifs et résultats}
\begin{columns}[T]
\begin{column}{0.49\textwidth}
\begin{tcolorbox}[title=Correctifs, colframe=primary, colback=primary!3]
\small
\textbf{Linguistiques :}
\begin{itemize}
    \item Restauration chaînes arabes
    \item Suppression prompt multilingue
    \item Stratégie deux passes
\end{itemize}
\textbf{Techniques :}
\begin{itemize}
    \item Listes fillers améliorées
    \item Détecteur acoustique nasaux (\texttt{librosa})
    \item Comptage répétitions
\end{itemize}
\end{tcolorbox}
\end{column}
\begin{column}{0.49\textwidth}
\begin{tcolorbox}[title=Résultats par langue, colframe=success, colback=success!3]
\small
\textbf{Anglais} \cmark~stable\\
\textbf{Français} \cmark~itérations nécessaires\\
\textbf{Arabe} \cmark~cas difficile
\end{tcolorbox}
\begin{tcolorbox}[title=Configuration retenue, colframe=accent, colback=accent!3]
\small
\texttt{medium + beam\_size 1}\\
Réduction 32--38\% temps total\\
Impact minimal sur précision
\end{tcolorbox}
\end{column}
\end{columns}
\end{frame}

\begin{frame}{Difficultés rencontrées et résolutions}
\begin{keybox}
\textbf{Problèmes techniques résolus pendant cette phase}
\end{keybox}
\vspace{0.35cm}
\begin{columns}[T]
\begin{column}{0.49\textwidth}
\begin{tcolorbox}[title=Faux "multi-axes", colframe=danger, colback=danger!3]
\small
\begin{itemize}
    \item Termes polysémiques faisaient croire à mélange voix/présence/scène
    \item \textbf{Solution :} détection restreinte à lexiques plus explicites
    \item Tests de non-régression en place
\end{itemize}
\end{tcolorbox}
\end{column}
\begin{column}{0.49\textwidth}
\begin{tcolorbox}[title=Cohérence profil--rapport, colframe=warning!60!black, colback=warning!8]
\small
\begin{itemize}
    \item Rapport écrit avant enrichissement LLM
    \item Risque de perdre préférences utilisateur
    \item \textbf{Solution :} serveur recharge profil si champ absent des métadonnées
\end{itemize}
\end{tcolorbox}
\end{column}
\end{columns}
\vspace{0.3cm}
\begin{tcolorbox}[colback=secondary!45,colframe=primary!25]
\small \textbf{Apprentissage :} chaque difficulté a été résolue avec une solution spécifique et testée
\end{tcolorbox}
\end{frame}

\begin{frame}{Analyse visuelle : Problème et solution}
\begin{columns}[T]
\begin{column}{0.49\textwidth}
\begin{tcolorbox}[title=Problème, colframe=danger, colback=danger!3]
\small
\textbf{Logique générique pénalisait vidéos smartphone}
\begin{itemize}
    \item Cas test révélateur : visage centré, mains visibles, regard caméra
    \item Résultats anormaux :
    \begin{itemize}
        \item Présence : 71\%
        \item Contact visuel : 10\%
        \item Mains visibles : 7\%
    \end{itemize}
    \item Cause : pipeline vision (MediaPipe + OpenCV) inadapté aux mobiles
\end{itemize}
\end{tcolorbox}
\end{column}
\begin{column}{0.49\textwidth}
\begin{tcolorbox}[title=Solution, colframe=success, colback=success!3]
\small
\textbf{Profils d'appareil :}
\begin{itemize}
    \item \texttt{laptop\_desktop} : réglage initial (pitch -32° à 8°, yaw ±15°)
    \item \texttt{tablet} : tolère décalages latéraux (yaw ±25°)
    \item \texttt{smartphone} : traitement spécifique (pitch -18° à 18°)
    \item \texttt{unknown} : auto-détection + tolérance générique
\end{itemize}
\textbf{Correctifs smartphone :} format portrait conservé, fréquence 2 fps (vs 1 fps), tolérances regard assouplies
\end{tcolorbox}
\end{column}
\end{columns}
\end{frame}

\begin{frame}{Analyse visuelle : Résultats et auto-détection}
\begin{columns}[T]
\begin{column}{0.49\textwidth}
\begin{tcolorbox}[title=Résultats smartphone, colframe=accent, colback=accent!3]
\small
\textbf{Amélioration spectaculaire sur cas test}
\begin{itemize}
    \item Présence : 71\% $\rightarrow$ \textbf{100\%} (+29\%)
    \item Contact visuel : 10\% $\rightarrow$ \textbf{80\%+} (+70\%)
    \item Mains visibles : 7\% $\rightarrow$ \textbf{60\%+} (+53\%)
\end{itemize}
\textbf{Observations :} laptop et tablette plus stables, ajustement supplémentaire pour tablette (usages paysage)
\end{tcolorbox}
\end{column}
\begin{column}{0.49\textwidth}
\begin{tcolorbox}[title=Auto-détection, colframe=primary, colback=primary!3]
\small
\textbf{Stratégie hiérarchique :}
\begin{itemize}
    \item 1. Métadonnées \texttt{ffprobe} + nom de fichier (tablet/phone keywords)
    \item 2. Rotation réelle + ratio d'aspect
    \item 3. Heuristique géométrique (portrait $\rightarrow$ smartphone, landscape $\rightarrow$ tablette/laptop)
\end{itemize}
\textbf{Résultats détection :}
\begin{itemize}
    \item Smartphone : auto-détecté via métadonnées/rotation
    \item Tablette : reconnue via métadonnées/nom de fichier
    \item Laptop/desktop : déduit par heuristique
\end{itemize}
\end{tcolorbox}
\end{column}
\end{columns}
\end{frame}

\begin{frame}{Stratégie LLM : Fine-tuning et transition Groq}
\begin{columns}[T]
\begin{column}{0.49\textwidth}
\begin{tcolorbox}[title=Fine-tuning engine-centric, colframe=warning!60!black, colback=warning!8]
\small
\textbf{Approche centrée sur le moteur :}
\begin{itemize}
    \item Données contraintes par sorties réelles du moteur (\texttt{calculate\_scores}, \texttt{generate\_recommendations})
    \item Enrichissements : niveau (Débutant/Intermédiaire/Avancé), objectif (Pitch/Soutenance/PFE/Entretien)
    \item Données générées par système (fidélité totale aux scores)
\end{itemize}
\textbf{Résultats :}
\begin{itemize}
    \item \xmark~Qwen 1.5B : résultats non convaincants
    \item \xmark~Limitations cohérence/généralisation sur long terme
    \item \cmark~Preuve de concept validée (Ollama local)
\end{itemize}
\end{tcolorbox}
\end{column}
\begin{column}{0.49\textwidth}
\begin{tcolorbox}[title=Transition Groq, colframe=success, colback=success!3]
\small
\textbf{Choix stratégique :} Llama 3.3 70B + Groq
\begin{itemize}
    \item \cmark~Rapidité exceptionnelle (LPU) : latence maîtrisée
    \item \cmark~Puissance de raisonnement : missions fines, sans répétitions
    \item \cmark~Quota gratuit : 100--120 sessions/jour
    \item \cmark~Fidélité JSON et rôle ``Master Coach''
\end{itemize}
\textbf{Avantages :} infrastructure cloud fiable pour usage produit, meilleure qualité que modèles compacts
\end{tcolorbox}
\end{column}
\end{columns}
\vspace{0.3cm}
\begin{tcolorbox}[colback=secondary!45,colframe=primary!25]
\small \textbf{Décision :} migration vers Groq pour production (V2) après validation fine-tuning local
\end{tcolorbox}
\end{frame}

\begin{frame}{Méthodologie fine-tuning : construction du jeu de données}
\begin{keybox}
\textbf{Jeu de données de référence (v2) construit en quatre phases méthodiques}
\end{keybox}
\vspace{0.35cm}
\begin{columns}[T]
\begin{column}{0.49\textwidth}
\begin{tcolorbox}[title=Phases 1-2 : Ancrage et couverture, colframe=primary, colback=primary!3]
\small
\textbf{Phase seed (données réelles) :}
\begin{itemize}
    \item Ancrage sur sessions réelles annotées manuellement
    \item Base réaliste pour l'entraînement
\end{itemize}
\textbf{Phase archétypes :}
\begin{itemize}
    \item Couverture spectre pédagogique complet
    \item Profils types : du scénario multi-failles à la haute maîtrise
    \item Échantillonnage stratifié des cas d'usage
\end{itemize}
\end{tcolorbox}
\end{column}
\begin{column}{0.49\textwidth}
\begin{tcolorbox}[title=Phases 3-4 : Sondage et distillation, colframe=success, colback=success!3]
\small
\textbf{Phase sondage (engine probing) :}
\begin{itemize}
    \item Échantillonnage stratifié du domaine des métriques
    \item Étiquetage par moteur comme fonction déterministe $y = f(x)$
    \item Couverture exhaustive des combinaisons possibles
\end{itemize}
\textbf{Phase distillation enseignant--élève :}
\begin{itemize}
    \item Modèle enseignant (Llama 3.3 via API) transforme sorties chiffrées en formulations exploitables
    \item Transfert de connaissances vers petit modèle cible (Qwen 1.5B)
\end{itemize}
\end{tcolorbox}
\end{column}
\end{columns}
\end{frame}

\begin{frame}{Architecture backend et sécurité}
\begin{columns}[T]
\begin{column}{0.49\textwidth}
\begin{tcolorbox}[title=Architecture hybride découplée, colframe=primary, colback=primary!3]
\small
\textbf{Flux de traitement :}
\begin{itemize}
    \item \textbf{Phase immédiate} : scoring déterministe (moteur audio/vision)
    \item \textbf{Phase différée} : enrichissement asynchrone (Worker Celery)
    \item \textbf{Suivi client} : polling statut jusqu'à disponibilité
\end{itemize}
\textbf{Technologies backend :}
\begin{itemize}
    \item FastAPI : API RESTful performante
    \item Celery : worker asynchrone pour enrichissement
    \item PostgreSQL : base de données relationnelle
    \item Redis : cache et file d'attente Celery
\end{itemize}
\end{tcolorbox}
\end{column}
\begin{column}{0.49\textwidth}
\begin{tcolorbox}[title=Sécurité et authentification, colframe=accent, colback=accent!3]
\small
\textbf{Flux SMTP professionnel :}
\begin{itemize}
    \item Envoi e-mail vérification à l'inscription
    \item Liens sécurisés réinitialisation mot de passe
    \item Pages dédiées Next.js (vérification/reset)
\end{itemize}
\textbf{Gestion de compte :}
\begin{itemize}
    \item Authentification JWT sécurisée
    \item Gestion autonome compte utilisateur
    \item Force mot de passe dynamique
\end{itemize}
\textbf{Garde-fous LLM :}
\begin{itemize}
    \item Filtre anti-tutoiement (regex)
    \item Rejet bilans vagues
    \item Fallback déterministe (0 rapport vide)
\end{itemize}
\end{tcolorbox}
\end{column}
\end{columns}
\end{frame}

\begin{frame}{Performance : Optimisations beam\_size=1}
\begin{columns}[T]
\begin{column}{0.48\textwidth}
\textbf{Comparaison avant/après (vidéo 5 min)}
\vspace{0.1cm}
\begin{tabular}{lcc}
\rowcolor{primary!20}\textbf{Métrique} & \textbf{Avant} & \textbf{Après} \\
\hline
Cold run & 12:12 & 8:17 \\
Hot run & 11:59 & 7:23 \\
\rowcolor{accent!20}\textbf{Réduction} & -- & \textbf{32--38\%} \\
RTF & ~2.25x & ~1.35x \\
\hline
\end{tabular}
\end{column}
\begin{column}{0.48\textwidth}
\begin{tcolorbox}[title=Impact beam\_size=1, colframe=accent, colback=accent!3]
\small
\textbf{Avantages :}
\begin{itemize}
    \item Réduction 39--43\% transcription
    \item Impact minimal sur précision
    \item Trade-off favorable (fillers/hésitations bien détectés)
\end{itemize}
\end{tcolorbox}
\vspace{0.2cm}
\begin{tcolorbox}[title=Cache et parallélisme, colframe=success, colback=success!3]
\small
\textbf{Optimisations :}
\begin{itemize}
    \item Cache : -100\% chargement modèles
    \item Cache : -89.5\% enrichissement
    \item Parallélisme : vision $\parallel$ transcription
    \item Scalabilité : temps linéaire (~1.35x)
\end{itemize}
\end{tcolorbox}
\end{column}
\end{columns}
\vspace{0.2cm}
\begin{keybox}
\textbf{Résultat :} vidéo 5 min traitée en 7:23 (hot run), acceptable pour usage produit
\end{keybox}
\end{frame}

\begin{frame}{Analyse de performance : Observations clés}
\begin{columns}[T]
\begin{column}{0.48\textwidth}
\textbf{Temps par phase (vidéo 5 min)}
\vspace{0.1cm}
\begin{tabular}{lcc}
\rowcolor{primary!20}\textbf{Phase} & \textbf{Cold run} & \textbf{Hot run} \\
\hline
Chargement modèles & 2:18 & 0:00 \\
Extraction audio & 0:16 & 0:16 \\
Transcription Whisper & 5:31 & 5:31 \\
Analyse vision & 0:17 & 0:17 \\
Enrichissement LLM & 0:55 & 0:19 \\
\rowcolor{accent!20}\textbf{TOTAL} & \textbf{8:17} & \textbf{7:23} \\
\hline
\end{tabular}
\end{column}
\begin{column}{0.48\textwidth}
\begin{tcolorbox}[title=Impact du cache, colframe=success, colback=success!3]
\small
\textbf{Réductions (cold $\rightarrow$ hot) :}
\begin{itemize}
    \item Chargement modèles : \textbf{-100\%}
    \item Enrichissement : \textbf{-89.5\%}
    \item Transcription : impact modéré (goulot principal)
    \item Total : \textbf{-10.9\%}
\end{itemize}
\textbf{Goulot d'étranglement :} Transcription = 72--93\% du temps
\end{tcolorbox}
\vspace{0.2cm}
\begin{tcolorbox}[title=Parallélisme, colframe=primary, colback=primary!3]
\small
\textbf{Exécution simultanée :}
\begin{itemize}
    \item Vision $\parallel$ transcription (gain parallèle)
    \item Enrichissement LLM très rapide (cache RAG + Groq)
    \item Scalabilité : temps linéaire (~1.35x temps réel)
\end{itemize}
\end{tcolorbox}
\end{column}
\end{columns}
\vspace{0.2cm}
\begin{tcolorbox}[colback=secondary!45,colframe=primary!25]
\small \textbf{Optimisations futures :} GPU (10--15x plus rapide), modèles plus légers (compromis qualité/vitesse), batch processing
\end{tcolorbox}
\end{frame}

\begin{frame}{Prochaines étapes et bilan}
\begin{columns}[T]
\begin{column}{0.49\textwidth}
\begin{tcolorbox}[title=Prochaines étapes, colframe=accent, colback=accent!3]
\small
\textbf{Priorité 1 : Accélération GPU}\\
\textit{Réduction 6--15x (7 min $\rightarrow$ 30--60 sec)}\\
\vspace{0.1cm}
\textbf{Priorité 2 : Intelligence Émotionnelle (EQ)}\\
\textit{Détection émotions : stress, confiance, enthousiasme}\\
\textit{MediaPipe Face Mesh + librosa}\\
\vspace{0.1cm}
\textbf{Priorité 3 : Interface multilingue}\\
\textit{Arabe, Français, Anglais (i18n + RTL)}\\
\vspace{0.1cm}
\textbf{Priorité 4 : Coaching en temps réel}\\
\textit{Streaming audio + WebRTC (transcription + analyse live)}
\end{tcolorbox}
\end{column}
\begin{column}{0.49\textwidth}
\begin{tcolorbox}[title=Bilan de cette phase, colframe=success, colback=success!3]
\small
\textbf{Accomplissements majeurs :}
\begin{itemize}
    \item \cmark~Coaching : templates $\rightarrow$ agent dynamique (Groq Llama 3.3)
    \item \cmark~ASR : fillers mieux détectés (3 langues, détecteur acoustique nasaux)
    \item \cmark~Vision : profils appareil + auto-détection hiérarchique
    \item \cmark~LLM : fine-tuning Qwen $\rightarrow$ Groq API (V2)
    \item \cmark~Performance : -32--38\% temps (beam\_size=1, cache)
    \item \cmark~UX/UI : professionnel, SMTP, exports PDF/MD/HTML
    \item \cmark~Backend : FastAPI + Celery + PostgreSQL + Redis
\end{itemize}
\end{tcolorbox}
\end{column}
\end{columns}
\vspace{0.3cm}
\begin{keybox}
\textbf{Résultat :} prototype technique $\rightarrow$ produit utilisable, performant et professionnel (vidéo 5 min traitée en 7:23)
\end{keybox}
\end{frame}

\setbeamertemplate{footline}{}
\begin{frame}[plain]
\begin{tikzpicture}[remember picture, overlay]
    \fill[primary] (current page.south west) rectangle (current page.north east);
    \node[text=white, align=center] at (current page.center) {
        {\Huge\bfseries Merci de votre attention}\\[0.45cm]
        {\Large Questions et démonstration}\\[0.9cm]
        \begin{tikzpicture}
            \fill[accent] (0,0) rectangle (4cm,0.14cm);
        \end{tikzpicture}
    };
\end{tikzpicture}
\end{frame}

\end{document}
